% Solutions of 14-dynamic-arrays.
\section*{Chapter \ref{ch:b1:dynamic-arrays} --- Dynamic Arrays and Amortised Analysis}

\begin{solution}{exo:b1:dynamic-arrays:1}
\omoutput[2-4]{code/bachelor-1/14-dynamic-arrays/outputs/exo_ops.txt}{The three lines.}
Six pushes leave \omterm{def:b1:dynamic-arrays:dynamic-array}{size} 6 in \omterm{def:b1:dynamic-arrays:dynamic-array}{capacity} 8, after copies $1 + 2 + 4 = 7$. The pops bring the \omterm{def:b1:dynamic-arrays:dynamic-array}{size} to 5, 4, 3, 2; at 2, $2 \leq
8/4$, and the \omterm{def:b1:dynamic-arrays:dynamic-array}{capacity} halves to 4 with 2 copies. Three pushes: \omterm{def:b1:dynamic-arrays:dynamic-array}{sizes} 3 and 4 fit, the fifth element doubles the
\omterm{def:b1:dynamic-arrays:dynamic-array}{capacity} to 8 with 4 copies: $7 + 2 + 4 = 13$.
\end{solution}

\begin{solution}{exo:b1:dynamic-arrays:2}
Doubling: resizes at pushes 2, 3, 5 and 9, copying $1 + 2 + 4 + 8 = 15$ elements. Adding 2 cells: capacities 1, 3, 5, 7,
9, 11, with resizes when the \omterm{def:b1:dynamic-arrays:dynamic-array}{size} is 1, 3, 5, 7 and 9: $1 + 3 + 5 + 7 + 9 = 25$ copies, already more, and growing
quadratically.
\end{solution}

\begin{solution}{exo:b1:dynamic-arrays:3}
After three puts: \code{[1 2 3 .]}, $h = 0$, $k = 3$. After two takes: $h = 2$, $k = 1$ (the queue is 3). Put 4 goes to
cell 3, put 5 to cell $(2 + 2) \bmod 4 = 0$, put 6 to cell 1: \code{[5 6 3 4]}, $h = 2$, $k = 4$, full. The last take
returns 3: $h = 3$, $k = 3$, the queue is 4, 5, 6.
\end{solution}

\begin{solution}{exo:b1:dynamic-arrays:4}
The increments from 0 to 8 flip 1, 2, 1, 3, 1, 2, 1 and 4 bits: 15 in all, $2 \cdot 8 - 1$.
\end{solution}

\begin{solution}{exo:b1:dynamic-arrays:5}
After four puts in four cells, \code{tail} has wrapped to 0, equal to \code{head}: the full buffer looks empty, and the
loop takes nothing. With two indices modulo $N$ there are only $N$ possible differences \code{tail - head} for $N + 1$
possible \omterm{def:b1:dynamic-arrays:dynamic-array}{sizes}, 0 to $N$. Two repairs: keep the count (as \cref{lst:b1:dynamic-arrays:ring} does), or never fill the last
cell, declaring the buffer full when \texttt{(tail + 1) \% N == head}, at the price of one unused cell.
\end{solution}

\begin{solution}{exo:b1:dynamic-arrays:6}
Half policy: at \omterm{def:b1:dynamic-arrays:dynamic-array}{size} 1024 in \omterm{def:b1:dynamic-arrays:dynamic-array}{capacity} 1024, a push doubles (1024 copies) to \omterm{def:b1:dynamic-arrays:dynamic-array}{size} 1025 in 2048; the pop brings the \omterm{def:b1:dynamic-arrays:dynamic-array}{size} to
$1024 \leq 2048/2$, and the \omterm{def:b1:dynamic-arrays:dynamic-array}{capacity} halves (1024 copies). Every operation copies 1024 elements. Quarter policy: the first
push copies 1024; then the \omterm{def:b1:dynamic-arrays:dynamic-array}{size} moves between 1024 and 1025 in \omterm{def:b1:dynamic-arrays:dynamic-array}{capacity} 2048, far from both 2048 and 512: no other copy,
so $1024/1000 \approx 1$ per operation. For $\Theta(m^2)$: $m/2 = 2^k$ pushes, then $m/2$ alternating operations of
$2^k = m/2$ copies each.
\end{solution}

\begin{solution}{exo:b1:dynamic-arrays:7}
\omcode{code/bachelor-1/14-dynamic-arrays/exercises/stack.c}{6}{25}{A stack on the \omterm{def:b1:dynamic-arrays:dynamic-array}{dynamic array}}
The header \code{exercises/stack.h} declares the six operations with \code{typedef struct stack stack;}. The top is the
last element of the \omterm{def:b1:pointers-and-arrays:array}{array}, so push and pop are those of the \omterm{def:b1:pointers-and-arrays:array}{array}, with \omterm{def:b1:dynamic-arrays:amortised}{amortised cost} $O(1)$. The test pushes and pops
1000 values and checks brackets: \code{([()[]])} is balanced, \code{([)]} and \code{((} are not.
\end{solution}

\begin{solution}{exo:b1:dynamic-arrays:8}
Bit $j$ changes at the increment from $x$ to $x + 1$ exactly when $2^j$ divides $x + 1$, which happens $\lfloor n/2^j
\rfloor$ times for $x + 1 = 1, \dots, n$. The total is $\sum_{j \geq 0} \lfloor n/2^j \rfloor < n \sum_j 2^{-j} = 2n$.
\end{solution}

\begin{solution}{exo:b1:dynamic-arrays:9}
Charge 2 credits per increment. An increment sets exactly one bit from 0 to 1 (unless it overflows): one credit pays that
flip, and the other is stored on the bit that is now 1. The flips from 1 to 0 are paid by the credits stored on those
bits. Every bit at 1 holds one credit, so the store is never negative, and $n$ increments flip at most $2n$ bits.
\end{solution}

\begin{solution}{exo:b1:dynamic-arrays:10}
Let $\Phi = \alpha s - c/(r - 1)$ for \omterm{def:b1:dynamic-arrays:dynamic-array}{size} $s$ and \omterm{def:b1:dynamic-arrays:dynamic-array}{capacity} $c$. A push without resize copies nothing and raises $\Phi$ by
$\alpha$. A push on a full \omterm{def:b1:pointers-and-arrays:array}{array} of \omterm{def:b1:dynamic-arrays:dynamic-array}{size} $s = c$ copies $c$ elements, and, the \omterm{def:b1:dynamic-arrays:dynamic-array}{capacity} growing by at least $(r - 1)c$,
$\Phi$ changes by at most $\alpha - (r - 1)c/(r - 1) = \alpha - c$: amortised $c + \alpha - c = \alpha$ copies. Right after a resize $s \geq c/r$, so $\Phi \geq c(\alpha/r -
1/(r - 1))$, which is non-negative for $\alpha = r/(r - 1)$; it stays so as pushes raise it. Hence $n$ pushes copy at most
$\frac{r}{r - 1} n$ plus a constant from $\Phi_0$ and the rounding of the new \omterm{def:b1:dynamic-arrays:dynamic-array}{capacity}. For $r = 1.5$ the bound is 3
(the table reaches 2.76), for $r = 2$ it is 2 (the table reaches 1.64).
\end{solution}

\begin{solution}{exo:b1:dynamic-arrays:11}
$\Phi \geq 0$ always, and $\Phi_0 = \frac12$ for the empty \omterm{def:b1:pointers-and-arrays:array}{array} of \omterm{def:b1:dynamic-arrays:dynamic-array}{capacity} 1. Pushes: in the upper half, $\Delta\Phi =
+2$; in the lower half, $-1$; on a full \omterm{def:b1:pointers-and-arrays:array}{array} of \omterm{def:b1:dynamic-arrays:dynamic-array}{size} $c$, $c$ copies and $\Phi$ goes from $c$ to 2, amortised $2 - c + c =
2$. Pops without resize: $-2$ in the upper half, $+1$ in the lower. A pop that shrinks leaves \omterm{def:b1:dynamic-arrays:dynamic-array}{size} $c/4$ in \omterm{def:b1:dynamic-arrays:dynamic-array}{capacity}
$c/2$: $\Phi$ goes from $c/2 - (c/4 + 1) = c/4 - 1$ to $2 \cdot c/4 - c/2 = 0$, with $c/4$ copies: amortised 1 (the
small capacities 4 and 2 give 1 and $\frac12$). Every \omterm{def:b1:dynamic-arrays:amortised}{amortised cost} is at most 2, so $m$ operations copy at most $2m +
\Phi_0 - \Phi_m \leq 2m + \frac12$, hence at most $2m$, an integer. The test checks it along 200\,000 random operations.
\end{solution}

\begin{solution}{exo:b1:dynamic-arrays:12}
\omcodes{code/bachelor-1/14-dynamic-arrays/exercises/gvec.c}{6-18,33-44}{The generic \omterm{def:b1:pointers-and-arrays:array}{array}: bytes and an element \omterm{def:b1:dynamic-arrays:dynamic-array}{size}}
Elements are copied with \code{memcpy}, \code{elem_size} bytes at a time, and the address of element $i$ is the start of
the block plus $i \cdot \text{elem\_size}$. The user \omterm{def:b1:c-types-and-control:cast}{casts} it back: \code{*(double *)gvec_at(d, 4000)}. The test pushes
5000 doubles and pops them back in reverse, then stores three \code{book} structures and changes one in place through
its address.
\end{solution}

\begin{solution}{pb:b1:dynamic-arrays:1}
\textbf{1.} The $t$ trailing 1s become 0 and the next bit, a 0, becomes 1: $t + 1$ flips.

\textbf{2.} Bits 0 to $t - 1$ are 1, bit $t$ is 0, the $K - t - 1$ others are free: $2^{K - t - 1}$ integers.

\textbf{3.} Bit $j$ flips from $x$ to $x + 1$ exactly when $2^j$ divides $x + 1$: $\lfloor n/2^j \rfloor$ times.

\textbf{4.} $\sum_j \lfloor n/2^j \rfloor < n \sum_j 2^{-j} = 2n$; for $n = 2^k$ the sum is $2^k + 2^{k-1} + \dots + 1 =
2^{k+1} - 1$.

\textbf{5.} 15, 127, 1023, 8191 and 65\,535 flips for $n = 8, 64, 512, 4096, 32\,768$: each is $2n - 1$.

\textbf{6.} From $2^K - 1$, all $K$ bits flip. Such increments are rare: one in $2^{K}$ from 0.

\textbf{7.} $2^{\lceil \log_2 n \rceil} - 1$: the resizes copy $1, 2, \dots, 2^j$ with $2^j < n$. For $n = 17$, 31.

\textbf{8.} A push pays its write and stores two credits on its element. When \omterm{def:b1:dynamic-arrays:dynamic-array}{capacity} $c$ fills, the $c/2$ elements
pushed since the last resize hold $c$ credits, one per copy. With 2 credits, only $c/2$ would be stored.

\textbf{9.} After a push the \omterm{def:b1:pointers-and-arrays:array}{array} holds more than half its \omterm{def:b1:dynamic-arrays:dynamic-array}{capacity}: either no resize happened and the \omterm{def:b1:dynamic-arrays:dynamic-array}{size} grew, or a
resize left \omterm{def:b1:dynamic-arrays:dynamic-array}{size} $c + 1$ in \omterm{def:b1:dynamic-arrays:dynamic-array}{capacity} $2c$. So $2s - c > 0$.

\textbf{10.} 3 for each of the seventeen pushes, as the dashed line of the figure shows.

\textbf{11.} The \omterm{def:b1:dynamic-arrays:dynamic-array}{capacity} is $2^{\lceil \log_2 n \rceil} < 2n$: the unused cells are fewer than the elements, nearly as
many right after a resize.

\textbf{12.} The \omterm{def:b1:dynamic-arrays:methods}{accounting method}: credits stored on elements show where the cost of a copy comes from.

\textbf{13.} \omterm{def:b1:dynamic-arrays:dynamic-array}{Size} $2^k$ in \omterm{def:b1:dynamic-arrays:dynamic-array}{capacity} $2^k$; the push copies $2^k$ elements into \omterm{def:b1:dynamic-arrays:dynamic-array}{capacity} $2^{k+1}$; the pop leaves \omterm{def:b1:dynamic-arrays:dynamic-array}{size}
$2^k$, half of it, and copies $2^k$ elements back into \omterm{def:b1:dynamic-arrays:dynamic-array}{capacity} $2^k$. The state is the same as after the fill.

\textbf{14.} Each operation then copies $2^k$ elements: $m \cdot 2^k$ copies.

\textbf{15.} A growth copies $c$ elements into \omterm{def:b1:dynamic-arrays:dynamic-array}{capacity} $2c$; a shrink copies $c/4$ elements into \omterm{def:b1:dynamic-arrays:dynamic-array}{capacity} $c/2$. Either
way the copied elements fill exactly half the new \omterm{def:b1:dynamic-arrays:dynamic-array}{capacity} $C$.

\textbf{16.} The next growth needs $C/2$ more pushes, the next shrink at least $C/4$ pops; a resize copies at most $C$
elements, so the copies are paid by at least $C/4$ operations: a first, crude bound of 4 per operation.

\textbf{17.} The worst sequences push up to just past a power of 2, sometimes after a shrink and a regrowth
(\texttt{+++-{}-{}+++}\dots). The largest ratio is $31/18 \approx 1.722$, at $m = 18$: seventeen pushes and one pop.

\textbf{18.} The resizes copy $1 + 2 + \dots + 2^k = 2^{k+1} - 1$ elements, and $\frac{2^{k+1} - 1}{2^k + 1} \to 2$.

\textbf{19.} The trailing 0s become 1 and the next 1 becomes 0. From $2^j$: $j + 1$ flips.

\textbf{20.} 16 per operation: from $2^{15} - 1$ (fifteen 1s), the increment flips 16 bits to reach $2^{15}$, the
decrement flips them back. The output shows 16\,000 flips for 1000 operations.

\textbf{21.} Each operation flips $\log_2 n + 1$ bits when the counter alternates between $2^k - 1$ and $2^k$ with $2^k
\approx n$: $\Theta(n \log n)$ in all.

\textbf{22.} Question 21 is about the standard representation, which is unique: $2^k - 1$ and $2^k$ differ in $k + 1$
bits, so any method that stores them must flip them. A redundant representation does not store either, but some
equivalent form of each, which leaves room to absorb a carry later.

\textbf{23.} Yes: the potential of \cref{exo:b1:dynamic-arrays:11} covers any mix. The quarter rule leaves slack in both
directions after each resize; the binary counter has no slack between $2^k - 1$ and $2^k$.

\textbf{24.} \emph{The elastic \omterm{def:b1:pointers-and-arrays:array}{array} theorem}: with growth $\times 2$ and shrinking to half when a quarter full, any $m$
pushes and pops from the empty \omterm{def:b1:pointers-and-arrays:array}{array} make at most $2m$ element copies (\cref{exo:b1:dynamic-arrays:11}). The constant 2 is
the smallest possible: $2^k + 1$ pushes make $2^{k+1} - 1$ copies, a ratio that approaches 2. The exhaustive search finds
at most 1.722 copies per operation for $m \leq 26$, at seventeen pushes and a pop.
\end{solution}
